\scanpage{167}
\uline{Variantes}\qquad a) On peut prendre sur $\deleted{X}$ un module de coeff. qu. coh. $E$, avec connexion intégrable rel. à $T$, d'où $\Omega^\cdot_{X/T}\otimes E$ avec structure de complexe. Tout ce qui précède s'applique, (grâce au fait que les suites exactes (4) splittent localement, donc \deleted{\unsure{restent exactes en tensorisant les filtrés|suite exacte en tensorisant les filtrations}} de même pour la filtration correspondante sur $\Omega^\cdot_{S/T}\otimes E$, donc
\[
 \mathrm{Gr}_{X/S}^{p}(\Omega^\cdot_{X/T}\otimes E)
 \simeq\mathrm{Gr}_{X/S}^{p}(\Omega^\cdot_{X/T})\otimes E
\]
du point de vue additif. Pour pouvoir appliquer la formule de projection, et obtenir, en analogie avec \inserted{(3)} \deleted{\unsure{(3)}},
\[
 \mathrm{Gr}_{X/S}^{p}(\mathbb R f_*(\Omega^\cdot_{X/T}\otimes E))
 \simeq\Omega^p_{S/T}\overset{L}{\otimes}\mathbb R f_*(\Omega^\cdot_{\unsure{Y|X}/S}\otimes E)[-p],
 \tag*{(21)}
\]
il semble qu'il faille supposer $E$ \uline{plat sur} $X$ \uline{ou} $\Omega^1_{S/T}$ \uline{plat sur} $S$. Sous cette réserve, on obtient des \deleted{stratifications} connexions intégrables sur les $\mathbb R f_*(\Omega^\cdot_{X/S}\otimes E)$ rel. à $T$, et une suite spectrale qui, dans le $2^{\text{ième}}$ cas ($\Omega^1_{S/T}$ plat) prend la forme simple
\[
 H^*_{\mathrm{DR}}(X/T,E)\Longleftarrow E_2^{pq}=H^p_{\mathrm{DR}}(S/T,\mathscr H^q_{\mathrm{DR}}(X/S,E))
 \tag*{(22)}
\]

b) On peut localiser la suite spectrale obtenue sur $T$, (plus généralement, relativement à un morphisme quelconque $S\longrightarrow R$). On peut dans ce cas \deleted{\illegible{1}}\inserted{formuler} des compatibilités de la suite spectrale obtenue sur $T$, avec des connexions intégrables relativement à une base en dessous de $T$.
% archive sheet number: 166

